\documentclass[11pt,a4paper]{article}
\usepackage{fontspec,amsmath,amssymb,amsthm,microtype}
\setmainfont{lmroman10-regular.otf}[BoldFont=lmroman10-bold.otf,ItalicFont=lmroman10-italic.otf,BoldItalicFont=lmroman10-bolditalic.otf]
\usepackage{xcolor}
\usepackage[most]{tcolorbox}
\usepackage{mathrsfs}
\providecommand{\lt}{<}
\providecommand{\gt}{>}
\setmonofont{DejaVuSansMono.ttf}[BoldFont=DejaVuSansMono-Bold.ttf,ItalicFont=DejaVuSansMono-Oblique.ttf,BoldItalicFont=DejaVuSansMono-BoldOblique.ttf,Scale=0.86]
\newfontfamily\leanfallback{STIXTwoMath-Regular.otf}[Scale=0.86]
\definecolor{leanrule}{gray}{0.62}
\definecolor{leanback}{gray}{0.962}
\definecolor{leantitle}{gray}{0.22}
\definecolor{leanmuted}{gray}{0.38}
\newtcolorbox{leanbox}[1]{enhanced,breakable,colback=leanback,colframe=leanrule,
  boxrule=0.4pt,arc=0.6mm,left=2.4mm,right=2.4mm,top=1.8mm,bottom=1.8mm,
  before skip=0.8em,after skip=1em,fonttitle=\small\bfseries,coltitle=leantitle,
  colbacktitle=leanback,titlerule=0.3pt,title={#1}}
\newlength{\leancharwidth}
\newenvironment{leancode}{\par\smallskip\begingroup\ttfamily\footnotesize\raggedright
  \settowidth{\leancharwidth}{M}\setlength{\parskip}{0pt}\setlength{\parindent}{0pt}
  \frenchspacing}{\par\endgroup\smallskip}
\newcommand{\leanline}[2]{\hangindent=\dimexpr#1\leancharwidth+4\leancharwidth\relax
  \hangafter=1\noindent\hspace*{#1\leancharwidth}#2\par}
\newcommand{\leansource}[1]{\par{\footnotesize\color{leanmuted}#1}\par}
\newcommand{\leanlabel}[1]{\par\smallskip\noindent{\small\bfseries #1}\par}
\usepackage[margin=25mm]{geometry}
\usepackage[unicode,hidelinks]{hyperref}
\hypersetup{pdfauthor={LangevinLean Project}}
\setlength{\emergencystretch}{3em}
\setlength{\parskip}{0.4em}
\setlength{\parindent}{0pt}
\allowdisplaybreaks[2]
\raggedbottom
\title{\texorpdfstring{Relative Fisher information and the frozen-gradient error}{Relative Fisher information and the frozen-gradient error}}
\author{LangevinLean Project}
\date{}
\begin{document}\maketitle
Draft mathematical exposition; these notes may contain mistakes.


Let \(\rho_t\) be the frozen Gaussian interpolation, \(r_t=d\rho_t/d\nu\), \(s_t=\nabla\log r_t\), and \(G_0=\int\|\nabla U\|^2d\mu<\infty\). The decomposition \(s_t=\nabla\log p_t+\nabla U\), the Gaussian score estimate, and the propagated gradient moment give

\[I_t\le4(1+L^2t^2)G_0+8tdL^2+d/t.\]

This estimate is uniformly finite on positive compact time intervals and proves local space-time integrability of \(r_t\|s_t\|^2\) under \(dt\,d\nu\). The same intervals have integrable frozen-gradient discrepancy, with the bound \(\mathbb E\|\nabla U(Y_t)-\nabla U(X)\|^2\le L^2(t^2G_0+2td)\). Changing variables to the initial/\allowbreak{}output pair law gives the identical bound with Gaussian kernel weight. These two energies justify the score pairings and removal of bounded-log tests.



Earlier results used here: \cite[Section~3.38]{LangevinLean2026ULASC}.

\section{\texorpdfstring{A square bound for the sum of two vectors}{A square bound for the sum of two vectors}}


The triangle inequality gives \(\|x+y\|\le\|x\|+\|y\|\). Square both nonnegative sides and use \(2\|x\|\|y\|\le\|x\|^2+\|y\|^2\), which follows from \((\|x\|-\|y\|)^2\ge0\). This proves \(\|x+y\|^2\le2\|x\|^2+2\|y\|^2\), the pointwise estimate used to split relative Fisher energy.



Expand \(\|a+b\|^2\), bound the mixed term by \(2\|a\|\|b\|\), and use \(2uv\le u^2+v^2\). This gives the factor two without requiring any relation between the vectors.

\begin{leanbox}{Lean: \texttt{norm\_\allowbreak{}add\_\allowbreak{}sq\_\allowbreak{}le\_\allowbreak{}two}}
\leansource{\texttt{norm\_\allowbreak{}add\_\allowbreak{}sq\_\allowbreak{}le\_\allowbreak{}two} \quad UnadjustedLangevin/\allowbreak{}LSI/\allowbreak{}GaussianEnergyIntegrability.lean, lines 28–32}
\begin{leancode}
\leanline{0}{private\ theorem\ norm\_add\_sq\_le\_two\ (x\ y\ :\ E)\ :}
\leanline{4}{‖x\ +\ y‖\ \textasciicircum{}\ 2\ ≤\ 2\ *\ ‖x‖\ \textasciicircum{}\ 2\ +\ 2\ *\ ‖y‖\ \textasciicircum{}\ 2}
\end{leancode}
\leanlabel{How the Lean proof follows the argument}
\texttt{norm\_\allowbreak{}add\_\allowbreak{}sq\_\allowbreak{}le\_\allowbreak{}two}: the intermediate facts include \texttt{h}, \texttt{hs}, in source order. The remaining scalar rearrangements are checked by \texttt{nlinarith}.
\end{leanbox}

\section{\texorpdfstring{Splitting the relative-score energy}{Splitting the relative-score energy}}


The jointly measurable score field agrees at positive time with the relative score. Its decomposition is \(s_t=\nabla U+\nabla\log p_t\). The propagated gradient is square integrable, and Gaussian smoothing gives \(\mathbb E_{\rho_t}\|\nabla\log p_t\|^2\le d/(2t)\). These facts also prove integrability of the relative-score square.



Apply the squared sum bound pointwise and integrate: \(I_t\le2G_t+2\mathbb E_{\rho_t}\|\nabla\log p_t\|^2\le2G_t+d/t\). The proof keeps the integrability witnesses for all terms before using integral monotonicity.



Split the relative score into the Lebesgue score and potential gradient. The preceding square inequality bounds its energy by twice each component energy, which are estimated by Gaussian smoothing and propagated score moments.

\begin{leanbox}{Lean: \texttt{interpolationRelativeScoreField\_\allowbreak{}energy\_\allowbreak{}le}}
\leansource{\texttt{interpolationRelativeScoreField\_\allowbreak{}energy\_\allowbreak{}le} \quad UnadjustedLangevin/\allowbreak{}LSI/\allowbreak{}GaussianEnergyIntegrability.lean, lines 36–68}
\begin{leancode}
\leanline{0}{theorem\ interpolationRelativeScoreField\_energy\_le}
\leanline{4}{(U\ :\ E\ →\ ℝ)\ (μ\ :\ Measure\ E)\ [IsProbabilityMeasure\ μ]}
\leanline{4}{[IsProbabilityMeasure\ (ULA.gibbsMeasure\ U)]}
\leanline{4}{(L\ :\ ℝ≥0)\ (hU\ :\ Differentiable\ ℝ\ U)\ (hLip\ :\ LipschitzWith\ L\ (gradient\ U))}
\leanline{4}{(hG\ :\ Integrable\ (fun\ x\ ↦\ ‖gradient\ U\ x‖\ \textasciicircum{}\ 2)\ μ)\ (t\ :\ ℝ)\ (ht\ :\ 0\ <\ t)\ :}
\leanline{4}{Integrable\ (fun\ y\ ↦\ ‖interpolationRelativeScoreField\ U\ μ\ t\ y‖\ \textasciicircum{}\ 2)}
\leanline{6}{(interpolationLaw\ U\ μ\ t)\ ∧}
\leanline{6}{(∫\ y,\ ‖interpolationRelativeScoreField\ U\ μ\ t\ y‖\ \textasciicircum{}\ 2\ ∂interpolationLaw\ U\ μ\ t)\ ≤}
\leanline{8}{2\ *\ (∫\ y,\ ‖gradient\ U\ y‖\ \textasciicircum{}\ 2\ ∂interpolationLaw\ U\ μ\ t)\ +}
\leanline{10}{(Module.finrank\ ℝ\ E\ :\ ℝ)\ /\ t}
\end{leancode}
\leanlabel{How the Lean proof follows the argument}
\texttt{interpolation\allowbreak{}Relative\allowbreak{}Score\allowbreak{}Field\_\allowbreak{}energy\_\allowbreak{}le}: the intermediate facts include \texttt{hgrad}, \texttt{hs}, \texttt{heq}, \texttt{hb}, \texttt{hdiv}, in source order; the source steps apply or rewrite \texttt{gradient\_\allowbreak{}log\_\allowbreak{}interpolation\allowbreak{}Relative\allowbreak{}Density}, \texttt{norm\_\allowbreak{}add\_\allowbreak{}sq\_\allowbreak{}le\_\allowbreak{}two}, \texttt{integral\_\allowbreak{}add}, \texttt{integral\_\allowbreak{}const\_\allowbreak{}mul}, \texttt{hb.\allowbreak{}trans}. The remaining scalar rearrangements are checked by \texttt{ring}, \texttt{nlinarith}.
\end{leanbox}


Earlier results used here: \cite[Section~3.26.3]{LangevinLean2026ULASC}.

\section{\texorpdfstring{An explicit finite relative-score energy bound}{An explicit finite relative-score energy bound}}


The interpolation gradient-moment estimate is \(G_t\le2(1+L^2t^2)G_0+4tdL^2\). Substitute it into \(I_t\le2G_t+d/t\) and distribute the factor two. This gives \(I_t\le4(1+L^2t^2)G_0+8tdL^2+d/t\). Every quantity on the right is determined by the initial gradient energy and the original potential; no entropy-production assumption is used.



Insert the Gaussian Fisher bound \(d/(2t)\) and the propagated gradient estimate. The result is finite at each positive time and uniformly integrable on intervals bounded away from zero.

\begin{leanbox}{Lean: \texttt{interpolationRelativeScoreField\_\allowbreak{}energy\_\allowbreak{}bound}}
\leansource{\texttt{interpolationRelativeScoreField\_\allowbreak{}energy\_\allowbreak{}bound} \quad UnadjustedLangevin/\allowbreak{}LSI/\allowbreak{}GaussianEnergyIntegrability.lean, lines 72–83}
\begin{leancode}
\leanline{0}{theorem\ interpolationRelativeScoreField\_energy\_bound}
\leanline{4}{(U\ :\ E\ →\ ℝ)\ (μ\ :\ Measure\ E)\ [IsProbabilityMeasure\ μ]}
\leanline{4}{[IsProbabilityMeasure\ (ULA.gibbsMeasure\ U)]}
\leanline{4}{(L\ :\ ℝ≥0)\ (hU\ :\ Differentiable\ ℝ\ U)\ (hLip\ :\ LipschitzWith\ L\ (gradient\ U))}
\leanline{4}{(hG\ :\ Integrable\ (fun\ x\ ↦\ ‖gradient\ U\ x‖\ \textasciicircum{}\ 2)\ μ)\ (t\ :\ ℝ)\ (ht\ :\ 0\ <\ t)\ :}
\leanline{4}{(∫\ y,\ ‖interpolationRelativeScoreField\ U\ μ\ t\ y‖\ \textasciicircum{}\ 2\ ∂interpolationLaw\ U\ μ\ t)\ ≤}
\leanline{6}{4\ *\ (1\ +\ (L\ :\ ℝ)\ \textasciicircum{}\ 2\ *\ t\ \textasciicircum{}\ 2)\ *\ (∫\ x,\ ‖gradient\ U\ x‖\ \textasciicircum{}\ 2\ ∂μ)\ +}
\leanline{8}{8\ *\ t\ *\ (Module.finrank\ ℝ\ E\ :\ ℝ)\ *\ (L\ :\ ℝ)\ \textasciicircum{}\ 2\ +}
\leanline{8}{(Module.finrank\ ℝ\ E\ :\ ℝ)\ /\ t}
\end{leancode}
\leanlabel{How the Lean proof follows the argument}
\texttt{interpolation\allowbreak{}Relative\allowbreak{}Score\allowbreak{}Field\_\allowbreak{}energy\_\allowbreak{}bound}: the intermediate fact is \texttt{h}, in source order. The remaining scalar rearrangements are checked by \texttt{nlinarith}.
\end{leanbox}


Earlier results used here: \cite[Section~3.26.3]{LangevinLean2026ULASC}.

\section{\texorpdfstring{The weighted density representation of score energy}{The weighted density representation of score energy}}


The explicitly constructed relative density satisfies \(r_t\nu=\rho_t\). The weighted-density integrability equivalence therefore transports integrability of \(\|s_t\|^2\) under \(\rho_t\) to integrability of \(r_t\|s_t\|^2\) under the fixed target \(\nu\). The corresponding integral identity is \(\int r_t\|s_t\|^2d\nu=\int\|s_t\|^2d\rho_t\). Strict positivity and measurability of the constructed density justify the ordinary weighted-integral formulas.



Changing from the proposed law to its density multiplies the score square by \(p_t\), or by \(r_t\) under Gibbs. These are equivalent representations of the same finite energy.

\begin{leanbox}{Lean: \texttt{interpolationRelativeScoreField\_\allowbreak{}weighted\_\allowbreak{}energy}}
\leansource{\texttt{interpolationRelativeScoreField\_\allowbreak{}weighted\_\allowbreak{}energy} \quad UnadjustedLangevin/\allowbreak{}LSI/\allowbreak{}GaussianEnergyIntegrability.lean, lines 87–108}
\begin{leancode}
\leanline{0}{theorem\ interpolationRelativeScoreField\_weighted\_energy}
\leanline{4}{(U\ :\ E\ →\ ℝ)\ (μ\ :\ Measure\ E)\ [IsProbabilityMeasure\ μ]}
\leanline{4}{[IsProbabilityMeasure\ (ULA.gibbsMeasure\ U)]}
\leanline{4}{(L\ :\ ℝ≥0)\ (hU\ :\ Differentiable\ ℝ\ U)\ (hLip\ :\ LipschitzWith\ L\ (gradient\ U))}
\leanline{4}{(hG\ :\ Integrable\ (fun\ x\ ↦\ ‖gradient\ U\ x‖\ \textasciicircum{}\ 2)\ μ)\ (t\ :\ ℝ)\ (ht\ :\ 0\ <\ t)\ :}
\leanline{4}{Integrable\ (fun\ y\ ↦\ interpolationRelativeDensity\ U\ μ\ t\ y\ *}
\leanline{6}{‖interpolationRelativeScoreField\ U\ μ\ t\ y‖\ \textasciicircum{}\ 2)\ (ULA.gibbsMeasure\ U)\ ∧}
\leanline{6}{(∫\ y,\ interpolationRelativeDensity\ U\ μ\ t\ y\ *}
\leanline{8}{‖interpolationRelativeScoreField\ U\ μ\ t\ y‖\ \textasciicircum{}\ 2\ ∂ULA.gibbsMeasure\ U)\ =}
\leanline{8}{∫\ y,\ ‖interpolationRelativeScoreField\ U\ μ\ t\ y‖\ \textasciicircum{}\ 2\ ∂interpolationLaw\ U\ μ\ t}
\end{leancode}
\leanlabel{How the Lean proof follows the argument}
\texttt{interpolation\allowbreak{}Relative\allowbreak{}Score\allowbreak{}Field\_\allowbreak{}weighted\_\allowbreak{}energy}: the intermediate facts include \texttt{hrm}, \texttt{hr}, \texttt{hi}, \texttt{hh}, in source order; the source steps apply or rewrite \texttt{tilted\allowbreak{}Measure\_\allowbreak{}interpolation\allowbreak{}Relative\allowbreak{}Density}, \texttt{ENNReal.\allowbreak{}to\allowbreak{}Real\_\allowbreak{}of\allowbreak{}Real}, \texttt{smul\_\allowbreak{}eq\_\allowbreak{}mul}.
\end{leanbox}


Earlier results used here: \cite[Section~3.26.3]{LangevinLean2026ULASC}.

\section{\texorpdfstring{Spacetime integrability of relative-score energy}{Spacetime integrability of relative-score energy}}


The relative density and the chosen relative-score field are jointly measurable in time and space. Their weighted energy is nonnegative. On a positive interval \([a,b]\), the section at every time is integrable by the density-change theorem.



Since \(0<a\le t\le b\), the quantitative estimate is bounded uniformly by \(4(1+L^2b^2)G_0+8bdL^2+d/a\). The integral of the section norm is exactly its nonnegative energy, so this finite scalar bound makes it integrable in time. The product criterion yields \(r_t\|s_t\|^2\in L^1(dt\,d\nu)\) on the interval. This is the spacetime dominating function required to pass from bounded entropy tests to the full relative Fisher term.



The explicit time bound is a sum of powers of time and \(1/t\). On a positive compact interval each is integrable, proving spacetime square integrability of the measurable relative-score field.

\begin{leanbox}{Lean: \texttt{integrable\_\allowbreak{}interpolationRelativeScoreField\_\allowbreak{}energy\_\allowbreak{}spacetime}}
\leansource{\texttt{integrable\_\allowbreak{}interpolationRelativeScoreField\_\allowbreak{}energy\_\allowbreak{}spacetime} \quad UnadjustedLangevin/\allowbreak{}LSI/\allowbreak{}GaussianEnergyIntegrability.lean, lines 112–149}
\begin{leancode}
\leanline{0}{theorem\ integrable\_interpolationRelativeScoreField\_energy\_spacetime}
\leanline{4}{(U\ :\ E\ →\ ℝ)\ (μ\ :\ Measure\ E)\ [IsProbabilityMeasure\ μ]}
\leanline{4}{[IsProbabilityMeasure\ (ULA.gibbsMeasure\ U)]}
\leanline{4}{(L\ :\ ℝ≥0)\ (hU\ :\ Differentiable\ ℝ\ U)\ (hLip\ :\ LipschitzWith\ L\ (gradient\ U))}
\leanline{4}{(hG\ :\ Integrable\ (fun\ x\ ↦\ ‖gradient\ U\ x‖\ \textasciicircum{}\ 2)\ μ)\ (a\ b\ :\ ℝ)\ (ha\ :\ 0\ <\ a)\ :}
\leanline{4}{Integrable\ (fun\ p\ :\ ℝ\ ×\ E\ ↦\ interpolationRelativeDensity\ U\ μ\ p.1\ p.2\ *}
\leanline{6}{‖interpolationRelativeScoreField\ U\ μ\ p.1\ p.2‖\ \textasciicircum{}\ 2)}
\leanline{6}{((volume.restrict\ (Icc\ a\ b)).prod\ (ULA.gibbsMeasure\ U))}
\end{leancode}
\leanlabel{How the Lean proof follows the argument}
\texttt{integrable\_\allowbreak{}interpolation\allowbreak{}Relative\allowbreak{}Score\allowbreak{}Field\_\allowbreak{}energy\_\allowbreak{}spacetime}: the intermediate facts include \texttt{G}, \texttt{C}, \texttt{h\allowbreak{}G0}, \texttt{hbpos}, \texttt{henergy}, in source order; the source steps apply or rewrite \texttt{interpolation\allowbreak{}Relative\allowbreak{}Density\_\allowbreak{}joint\_\allowbreak{}measurable}, \texttt{strongly\allowbreak{}Measurable.\allowbreak{}mul}, \texttt{henergy.\allowbreak{}trans}, \texttt{ht.\allowbreak{}1}, \texttt{ht.\allowbreak{}2}.
\end{leanbox}


Earlier results used here: \cite[Section~3.26.3]{LangevinLean2026ULASC}.

\section{\texorpdfstring{The frozen-gradient discrepancy on the independent-noise space}{The frozen-gradient discrepancy on the independent-noise space}}


For the fixed independent pair \((X,Z)\sim\mu\otimes\mathcal N(0,I)\), define \(e_t(X,Z)=\nabla U(X-t\nabla U(X)+\sqrt{2t}Z)-\nabla U(X)\). The underlying probability space is independent of time. This avoids assuming a jointly measurable choice of a state-space conditional expectation when integrating the error energy over time.



The discrepancy is evaluated at the independent-noise proposal and its initial state. Its measurability follows from the continuous gradient and update map, fixing one field for subsequent pair integrals.

\begin{leanbox}{Lean: \texttt{frozenGradientErrorField}}
\leansource{\texttt{frozenGradientErrorField} \quad UnadjustedLangevin/\allowbreak{}LSI/\allowbreak{}GaussianEnergyIntegrability.lean, lines 152–153}
\begin{leancode}
\leanline{0}{def\ frozenGradientErrorField\ (U\ :\ E\ →\ ℝ)\ (t\ :\ ℝ)\ (p\ :\ E\ ×\ E)\ :\ E\ :=}
\leanline{2}{gradient\ U\ (ULA.malaProposal\ U\ t\ p.1\ p.2)\ -\ gradient\ U\ p.1}
\end{leancode}
\leanlabel{How the Lean proof follows the argument}
\texttt{frozen\allowbreak{}Gradient\allowbreak{}Error\allowbreak{}Field} is a definition with defining expression \texttt{gradient U (ULA.\allowbreak{}mala\allowbreak{}Proposal U t p.\allowbreak{}1 p.\allowbreak{}2) -\allowbreak{} gradient U p.\allowbreak{}1}. The right-hand side specifies the object used in the mathematical statement.
\end{leanbox}


Earlier results used here: \cite[Section~3.34.1]{LangevinLean2026ULASC}.

\section{\texorpdfstring{Spacetime integrability of discrepancy energy}{Spacetime integrability of discrepancy energy}}


The proposal and Lipschitz gradient make \((t,X,Z)\mapsto\|e_t(X,Z)\|^2\) jointly measurable. For \(t\ge0\), the exact Gaussian calculation already gives integrability and \(\mathbb E\|e_t\|^2\le L^2(t^2G_0+2td)\). On a compact nonnegative time interval this is at most \(L^2(b^2G_0+2bd)\). The time measure is finite, so Fubini's integrability criterion proves the squared error is integrable on the time/\allowbreak{}initial-state/\allowbreak{}noise product.



The pointwise gradient Lipschitz estimate bounds its square by \(L^2\) times displacement square. The exact displacement formula is polynomial in time with finite coefficients, giving an integrable spacetime energy.

\begin{leanbox}{Lean: \texttt{integrable\_\allowbreak{}frozenGradientErrorField\_\allowbreak{}energy\_\allowbreak{}spacetime}}
\leansource{\texttt{integrable\_\allowbreak{}frozenGradientErrorField\_\allowbreak{}energy\_\allowbreak{}spacetime} \quad UnadjustedLangevin/\allowbreak{}LSI/\allowbreak{}GaussianEnergyIntegrability.lean, lines 157–192}
\begin{leancode}
\leanline{0}{theorem\ integrable\_frozenGradientErrorField\_energy\_spacetime}
\leanline{4}{(U\ :\ E\ →\ ℝ)\ (μ\ :\ Measure\ E)\ [IsProbabilityMeasure\ μ]}
\leanline{4}{(L\ :\ ℝ≥0)\ (hLip\ :\ LipschitzWith\ L\ (gradient\ U))}
\leanline{4}{(hG\ :\ Integrable\ (fun\ x\ ↦\ ‖gradient\ U\ x‖\ \textasciicircum{}\ 2)\ μ)\ (a\ b\ :\ ℝ)\ (ha\ :\ 0\ ≤\ a)\ :}
\leanline{4}{Integrable\ (fun\ p\ :\ ℝ\ ×\ (E\ ×\ E)\ ↦\ ‖frozenGradientErrorField\ U\ p.1\ p.2‖\ \textasciicircum{}\ 2)}
\leanline{6}{((volume.restrict\ (Icc\ a\ b)).prod\ (μ.prod\ (stdGaussian\ E)))}
\end{leancode}
\leanlabel{How the Lean proof follows the argument}
\texttt{integrable\_\allowbreak{}frozen\allowbreak{}Gradient\allowbreak{}Error\allowbreak{}Field\_\allowbreak{}energy\_\allowbreak{}spacetime}: the intermediate facts include \texttt{C}, \texttt{hg}, \texttt{hm}, \texttt{ht0}, \texttt{h\allowbreak{}G0}, in source order; the source steps apply or rewrite \texttt{integrable\_\allowbreak{}prod\_\allowbreak{}iff}, \texttt{hm.\allowbreak{}aestrongly\allowbreak{}Measurable}, \texttt{integral\_\allowbreak{}nonneg}, \texttt{h.\allowbreak{}trans}, \texttt{ht.\allowbreak{}2}. The \texttt{calc} block checks the successive equalities or inequalities in the argument. The remaining scalar rearrangements are checked by \texttt{ring}.
\end{leanbox}


Earlier results used here: \cite[Section~3.34.1]{LangevinLean2026ULASC}.

\section{\texorpdfstring{The gradient-error energy under the Gaussian pair law}{The gradient-error energy under the Gaussian pair law}}


Map the independent initial/\allowbreak{}noise pair by \(T_t(x,z)=(x,x-t\nabla U(x)+\sqrt{2t}z)\). The squared error as a function of the output pair is \(F(x,y)=\|\nabla U(y)-\nabla U(x)\|^2\). Its composition with \(T_t\) is exactly the frozen-gradient discrepancy from the Gaussian energy calculation. Measurability of \(T_t\) and continuity of \(F\) let the map-measure integrability and integral theorems transport both the integrability witness and the bound \(L^2(t^2G_0+2td)\) to the initial/\allowbreak{}output pair law.



The pair-law change of variables transports the discrepancy energy from independent noise to the Gaussian starting-observation pair. Because the integrand is nonnegative, the identity can be established before signed Fubini is needed.

\begin{leanbox}{Lean: \texttt{frozenGaussianPair\_\allowbreak{}gradient\_\allowbreak{}error\_\allowbreak{}energy}}
\leansource{\texttt{frozenGaussianPair\_\allowbreak{}gradient\_\allowbreak{}error\_\allowbreak{}energy} \quad UnadjustedLangevin/\allowbreak{}LSI/\allowbreak{}GaussianEnergyIntegrability.lean, lines 197–217}
\begin{leancode}
\leanline{0}{theorem\ frozenGaussianPair\_gradient\_error\_energy}
\leanline{4}{(U\ :\ E\ →\ ℝ)\ (μ\ :\ Measure\ E)\ [IsProbabilityMeasure\ μ]}
\leanline{4}{(L\ :\ ℝ≥0)\ (hLip\ :\ LipschitzWith\ L\ (gradient\ U))}
\leanline{4}{(hG\ :\ Integrable\ (fun\ x\ ↦\ ‖gradient\ U\ x‖\ \textasciicircum{}\ 2)\ μ)\ (t\ :\ ℝ)\ (ht\ :\ 0\ ≤\ t)\ :}
\leanline{4}{let\ P\ :=\ (μ.prod\ (stdGaussian\ E)).map}
\leanline{6}{(fun\ z\ :\ E\ ×\ E\ ↦\ (z.1,\ ULA.malaProposal\ U\ t\ z.1\ z.2))}
\leanline{4}{Integrable\ (fun\ z\ :\ E\ ×\ E\ ↦\ ‖gradient\ U\ z.2\ -\ gradient\ U\ z.1‖\ \textasciicircum{}\ 2)\ P\ ∧}
\leanline{6}{(∫\ z\ :\ E\ ×\ E,\ ‖gradient\ U\ z.2\ -\ gradient\ U\ z.1‖\ \textasciicircum{}\ 2\ ∂P)\ ≤}
\leanline{8}{(L\ :\ ℝ)\ \textasciicircum{}\ 2\ *\ (t\ \textasciicircum{}\ 2\ *\ (∫\ x,\ ‖gradient\ U\ x‖\ \textasciicircum{}\ 2\ ∂μ)\ +}
\leanline{10}{2\ *\ t\ *\ (Module.finrank\ ℝ\ E\ :\ ℝ))}
\end{leancode}
\leanlabel{How the Lean proof follows the argument}
\texttt{frozen\allowbreak{}Gaussian\allowbreak{}Pair\_\allowbreak{}gradient\_\allowbreak{}error\_\allowbreak{}energy}: the intermediate facts include \texttt{h\allowbreak{}Y}, \texttt{h\allowbreak{}T}, \texttt{hg}, \texttt{h\allowbreak{}F}, \texttt{herror}, in source order; the source steps apply or rewrite \texttt{integrable\_\allowbreak{}map\_\allowbreak{}measure}, \texttt{h\allowbreak{}F.\allowbreak{}aestrongly\allowbreak{}Measurable}, \texttt{herror.\allowbreak{}1}, \texttt{integral\_\allowbreak{}map\_\allowbreak{}of\_\allowbreak{}strongly\allowbreak{}Measurable}, \texttt{herror.\allowbreak{}2}.
\end{leanbox}


Earlier results used here: \cite[Section~3.34.1]{LangevinLean2026ULASC}; \cite[Section~3.36.5]{LangevinLean2026ULASC}.

\section{\texorpdfstring{The kernel representation of gradient-error energy}{The kernel representation of gradient-error energy}}


The pair law is identified with \(K_t(x,y)\,\mu(dx)dy\) by the previously proved Gaussian proposal pair-density theorem. The proof specializes its measurable center to \(x-t\nabla U(x)\); the precision and step-size versions of the kernel agree pointwise.



Rewrite the pair-law energy theorem using this exact density identity. The weighted-density integrability equivalence gives \(K_t\|\nabla U(y)-\nabla U(x)\|^2\in L^1(\mu\otimes dy)\), and its integral formula gives the sharp bound \(\int K_t\|\nabla U(y)-\nabla U(x)\|^2d\mu\,dy\le L^2(t^2G_0+2td)\). The change of measure preserves the constants in this error-energy bound, which are used in the entropy recurrence.



Rewriting the pair measure using its conditional Gaussian density gives the kernel-weighted energy. Thus the same sharp bound is available in the weak PDE's Lebesgue integral representation.

\begin{leanbox}{Lean: \texttt{frozenGaussianKernel\_\allowbreak{}gradient\_\allowbreak{}error\_\allowbreak{}energy}}
\leansource{\texttt{frozenGaussianKernel\_\allowbreak{}gradient\_\allowbreak{}error\_\allowbreak{}energy} \quad UnadjustedLangevin/\allowbreak{}LSI/\allowbreak{}GaussianEnergyIntegrability.lean, lines 222–263}
\begin{leancode}
\leanline{0}{theorem\ frozenGaussianKernel\_gradient\_error\_energy}
\leanline{4}{(U\ :\ E\ →\ ℝ)\ (μ\ :\ Measure\ E)\ [IsProbabilityMeasure\ μ]}
\leanline{4}{(L\ :\ ℝ≥0)\ (hLip\ :\ LipschitzWith\ L\ (gradient\ U))}
\leanline{4}{(hG\ :\ Integrable\ (fun\ x\ ↦\ ‖gradient\ U\ x‖\ \textasciicircum{}\ 2)\ μ)\ (t\ :\ ℝ)\ (ht\ :\ 0\ <\ t)\ :}
\leanline{4}{Integrable\ (fun\ z\ :\ E\ ×\ E\ ↦\ realGaussianKernel\ (1\ /\ (4\ *\ t))}
\leanline{6}{(z.1\ -\ t\ •\ gradient\ U\ z.1)\ z.2\ *\ ‖gradient\ U\ z.2\ -\ gradient\ U\ z.1‖\ \textasciicircum{}\ 2)}
\leanline{6}{(μ.prod\ volume)\ ∧}
\leanline{6}{(∫\ z\ :\ E\ ×\ E,\ realGaussianKernel\ (1\ /\ (4\ *\ t))}
\leanline{8}{(z.1\ -\ t\ •\ gradient\ U\ z.1)\ z.2\ *\ ‖gradient\ U\ z.2\ -\ gradient\ U\ z.1‖\ \textasciicircum{}\ 2}
\leanline{10}{∂μ.prod\ volume)\ ≤}
\leanline{8}{(L\ :\ ℝ)\ \textasciicircum{}\ 2\ *\ (t\ \textasciicircum{}\ 2\ *\ (∫\ x,\ ‖gradient\ U\ x‖\ \textasciicircum{}\ 2\ ∂μ)\ +}
\leanline{10}{2\ *\ t\ *\ (Module.finrank\ ℝ\ E\ :\ ℝ))}
\end{leancode}
\leanlabel{How the Lean proof follows the argument}
\texttt{frozen\_\allowbreak{}displacement\_\allowbreak{}second\_\allowbreak{}moment} expands the squared displacement of the product-space Gaussian update. The mixed expectation vanishes using independence and centering; the standard Gaussian second moment supplies the dimensional term. \texttt{frozen\allowbreak{}Gaussian\allowbreak{}Kernel\_\allowbreak{}gradient\_\allowbreak{}error\_\allowbreak{}energy} transports that calculation to the kernel input-output measure, then applies the Lipschitz gradient bound to the squared gradient difference. The relevant integrability lemmas precede the signed integral expansion. The formulas retain the initial gradient moment and do not introduce a position moment.
\end{leanbox}


Earlier results used here: \cite[Section~3.34.1]{LangevinLean2026ULASC}; \cite[Section~3.37.12]{LangevinLean2026ULASC}; \cite[Section~3.37.13]{LangevinLean2026ULASC}; \cite[Section~3.37.4]{LangevinLean2026ULASC}; \cite[Section~3.38.1]{LangevinLean2026ULASC}.

\begin{thebibliography}{ULA-LSI}
\bibitem[ULA-SC]{LangevinLean2026ULASC}
LangevinLean Project. \emph{Unadjusted Langevin under Strong Convexity}. LL-01, 2026. Draft report, content revision d34ead06987c. \url{https://odakyildiz.github.io/LangevinLean/docs/pdf/UnadjustedLangevin/FullProof.pdf}.
\end{thebibliography}

\end{document}
